\subsection{Overview and Reconstruction Formulation}
\label{sec:method_overview}

This work develops a supervised \gls{ml} pipeline for separating overlapping
electron showers in the \gls{ldmx} \gls{ecal}. The method starts from
reconstructed detector collections stored in simulated \gls{ldmxsw} ROOT
events and produces hit-level assignments, together with optional
event-level electron-count predictions. It is therefore a detector-specific
reconstruction study: it is motivated by the global, object-based view of
\gls{mlpf} \cite{pata2021-mlpf,pata2026-mlpf}, but it does not attempt to
construct the complete particle-flow output used in a collider experiment.

For one event, the observed \gls{ecal} response is represented as the
variable-size set
\begin{equation}
    \mathcal{H} =
    \left\{
    \left(\boldsymbol{r}_i,E_i\right)
    \right\}_{i=1}^{N_{\mathrm{ECal}}},
    \qquad
    \boldsymbol{r}_i=(x_i,y_i,z_i),
\end{equation}
where \(E_i\) is the reconstructed energy of hit \(i\). Depending on the
model, this set is augmented with a second variable-size set
\(\mathcal{T}\) of trigger-pad-track observations. The principal supervised
task is to assign each selected \gls{ecal} hit to the electron that gives its
dominant simulated deposited-energy contribution. An extended task also
predicts the composition of shared hits and whether an event contains two or
three electrons.

This representation preserves individual detector objects rather than
projecting the event to an image. The number of hits is event dependent, the
\gls{ecal} is three-dimensional, and overlapping showers need not be aligned
with a fixed rasterization. Moreover, trigger-pad-track information is not an
\gls{ecal} pixel quantity. Sets, learned graphs and attention-based token
sequences provide natural ways of relating these heterogeneous objects while
remaining insensitive to an arbitrary ordering of the input rows.

\subsection{Data and Simulation Inputs}
\label{sec:method_data}

\subsubsection{ROOT Collections}

The input events are simulated with the \gls{geant4}-based \gls{ldmxsw}
framework and stored in a ROOT tree named \texttt{LDMX\_Events}. The
supervised pipeline reads reconstructed \gls{ecal} hits from
\texttt{EcalRecHits\_overlay}. For each reconstructed hit it extracts the
cell identifier, Cartesian position, reconstructed energy and the stored
noise flag. Truth association is obtained from
\texttt{EcalSimHitsOverlay\_overlay}, which contains, per detector cell, the
track identifiers, energy depositions and origin identifiers of the
contributing simulated particles. Reconstructed and simulated hit records are
aligned by detector cell identifier. Thus, the network inputs remain
reconstructed quantities, whereas simulation information is used to construct
training targets.

The context-aware models additionally read
\texttt{TriggerPadTracks\_overlay}. In the implemented pipeline, each such
object contributes its \texttt{centroid} coordinate and its number of
photoelectrons, \texttt{pe}. These objects are treated as context nodes or
tokens: they can influence an \gls{ecal} prediction, but do not receive
hit-origin supervision themselves.

\subsubsection{Overlay Event Samples}

The data-loading and training code supports separate overlay sources marked
by their electron multiplicity. In particular, it accepts two-electron
(\(2e\)) and three-electron (\(3e\)) sources and can combine them in one
processed dataset. The fixed-three-origin development sample is useful for
testing shower assignment, while mixed \(2e/3e\) data enable the event-level
electron-count task. Since the conclusions drawn from a trained model depend
on the particular selected subset and split, sample sizes and numerical
performance are reported with the corresponding results rather than assumed
in the definition of the method.

\subsection{Supervised Targets}
\label{sec:method_targets}

\subsubsection{Dominant-Contribution Hit Labels}

A reconstructed hit may contain simulated energy deposited by more than one
overlaid electron. If \(e_{ij}\) is the energy contribution to reconstructed
hit \(i\) associated with origin \(j\), its hard assignment target is defined
by
\begin{equation}
    y_i^{\mathrm{origin}} =
    \argmax_{j \in \{1,2,3\}} e_{ij}.
    \label{eq:dominant_origin_target}
\end{equation}
The implemented ROOT-to-tensor conversion selects the origin identifier of
the largest deposited-energy contribution and maps accepted physical origin
identifiers \(1,2,3\) to class indices appropriate for the training loss.
This target asks which electron dominates the reconstructed hit; it does not
assert that a hit with contributions from several showers is physically
unmixed.

For the multi-task models, the deposited-energy contributions also define a
soft target,
\begin{equation}
    f_{ij} =
    \frac{e_{ij}}{\sum_k e_{ik}},
    \qquad
    \sum_j f_{ij}=1
\end{equation}
when all retained contributions correspond to the configured electron
origins. Predicting \(f_{ij}\) retains information about hits near an overlap
region that is discarded by the dominant-origin label alone.

\subsubsection{Noise and Canonical Electron Slots}

Hits marked as detector noise are excluded from the maintained hit-origin
baseline comparisons. The advanced slot-model path can instead retain these
hits and assign them to an explicit background class \(0\), with a
background-only fraction target. Noise-labelled hits are not used to
determine the presence or ordering of electron slots.

An origin identifier introduced during overlay is a provenance label rather
than an ordered physical observable: exchanging the identifiers of two
electrons leaves the reconstruction problem unchanged. To remove this
permutation ambiguity, the maintained comparisons use a canonical spatial
target. For each electron present in an event, the mean \(y\)-coordinate of
its selected \gls{ecal} hits is calculated. Electrons are ordered by this
mean coordinate, and their targets are remapped to canonical slots from
lowest to highest \(y\). The same reordering is applied to the fraction
targets. The original origin identifiers are retained in the processed
events for provenance. Other spatial axes are implemented for controlled
studies, while canonical ordering in \(y\) is the standard comparison
configuration because it provides a consistent relation to the available
one-dimensional trigger-pad centroid information.

\subsection{Feature Construction and Event Representation}
\label{sec:method_features}

The ECal-only representation supplies one row
\((x_i,y_i,z_i,E_i)\) for each retained reconstructed hit. The combined
representation places \gls{ecal} hits and trigger-pad-track objects in one
event tensor with the feature layout shown in
Table~\ref{tab:combined_features}. Type-indicator columns identify the
detector collection, while non-applicable continuous fields are zero-filled.
Boolean masks retain which rows correspond to ECal hits and which correspond
to trigger-pad tracks.

\begin{table}[H]
    \centering
    \caption{Implemented feature layout for the combined ECal and
    trigger-pad-track representation.}
    \label{tab:combined_features}
    \begin{tabular}{lll}
        \hline
        Columns & ECal hit row & Trigger-pad-track row \\
        \hline
        \texttt{is\_ecal}, \texttt{is\_tpad} & \(1,0\) & \(0,1\) \\
        \texttt{ecal\_x}, \texttt{ecal\_y}, \texttt{ecal\_z} &
            \(x,y,z\) & \(0,0,0\) \\
        \texttt{ecal\_energy} & \(E\) & \(0\) \\
        \texttt{tpad\_centroid}, \texttt{tpad\_pe} &
            \(0,0\) & \(c,n_{\mathrm{pe}}\) \\
        \hline
    \end{tabular}
\end{table}

All continuous columns in this canonical combined representation are
standardized using means and standard deviations calculated only from the
training split. The fitted statistics are then applied to validation and test
events and stored with the run artifacts. The detector-type indicator columns
are not standardized.

The maintained models process one variable-length event at a time. Several
events are accumulated before one optimizer update, but no fixed maximum
number of detector objects is imposed by the main training pipeline. For
models receiving both detector collections, the \texttt{ecal\_mask} restricts
the hit-level loss and hit-level metrics to \gls{ecal} rows. This distinction
allows trigger-pad-track tokens to provide information without requiring an
artificial hit-origin target for them.

\subsection{Dataset Preparation and Storage}
\label{sec:method_processing}

ROOT reading is implemented separately from tensorization and model training.
The reader uses \texttt{uproot} and jagged arrays to extract the required
collections in chunks, aligns truth contributions to reconstructed hits, and
passes plain event records to the tensorization step. Tensorization constructs
features, masks, dominant labels and, when requested, fraction and noise
targets. This separation makes the physics definition of a target independent
of the particular neural-network architecture.

For repeated training, tensorized events can be written as PyTorch
\texttt{.pt} data rather than being reconstructed from ROOT for every run.
Two storage forms are implemented: per-event cached tensors for configurable
subsets and a scalable sharded cache in which one shard corresponds to one
source ROOT file. A manifest records the source inputs, feature layout,
allowed labels and noise policy; the sharded form additionally stores an
index of event ranges and supports lazy loading with a small in-memory shard
cache. This design reduces repeated ROOT input/output and permits model
comparisons to use the same processed event definition.

Events are randomly assigned, with a fixed seed, to training, validation and
test subsets in proportions \(80\%\), \(15\%\) and \(5\%\), respectively.
The normalization statistics, optional class weights and parameter updates
are determined from the training subset only. Validation data are used for
checkpoint selection and optional learning-rate scheduling, while test data
are held for final evaluation.

\subsection{Model Architectures}
\label{sec:method_models}

\subsubsection{Hit-Origin Baseline Models}

Four maintained classifiers establish controlled comparisons of input
information and interaction mechanism: ECal-only and ECal-plus-trigger-pad
versions of a \gls{gravnet} model, and the corresponding two versions of a
transformer encoder. All four receive the same canonical targets and return
one vector of origin-class logits for each input ECal hit. In a
context-aware version, predictions emitted for trigger-pad-track rows are
discarded by the ECal supervision mask.

In the GravNet classifiers, an input multilayer perceptron first maps each
node to a latent representation. Each \texttt{GravNetConv} layer learns
coordinates in a latent space, connects neighbouring nodes in that learned
space and propagates features between them. Residual updates preserve the
previous representation before a linear classification head produces
hit-origin logits. This use of learned neighbourhoods is suited to
high-granularity calorimeter showers whose useful adjacency is not fully
specified by Euclidean detector distance alone \cite{qasim2019-gravnet}.

In the transformer classifiers, each hit or mixed-detector token is embedded
by an input multilayer perceptron and passed through a stack of encoder
layers. Full self-attention permits every object in one event to condition its
representation on every other object. A linear per-token head then predicts
the canonical dominant-origin class. This architecture follows the
set-to-set aspect of recent \gls{mlpf} reconstruction methods while targeting
the narrower LDMX shower-separation problem \cite{pata2026-mlpf}.

\subsubsection{Multi-Task MLPF-Lite Model}

The implemented MLPF-lite model uses the combined ECal and trigger-pad-track
tensor and a shared transformer encoder. It has two ECal-supervised heads:
one predicts dominant-origin logits and the other predicts the soft
origin-energy fractions. It is consequently able to study whether a shared
event representation can describe both a discrete shower assignment and the
degree of local shower mixing. This model is MLPF-inspired in its
heterogeneous token representation and multi-task prediction, but its
outputs remain LDMX-specific ECal assignments rather than complete
particle-flow objects.

\subsubsection{Slot and Electron-Count Model}

The slot model extends the same context-token transformer to the mixed
\(2e/3e\) setting. Its hit-level outputs consist of a background class and up
to three canonical electron slots, together with corresponding fraction
predictions. To obtain event-level outputs, encoded token features are pooled
using their mean, maximum and the logarithm of the token multiplicity. The
resulting event representation is passed to heads that predict the validity
of each electron slot and the number of electrons in the event. The model can
therefore connect local shower assignment with the event-level question of
whether an additional overlapping electron is present.

\subsection{Training Objective and Optimization}
\label{sec:method_training}

For the baseline hit-origin classifiers, the loss is categorical
cross-entropy evaluated only on retained \gls{ecal} hits,
\begin{equation}
    \mathcal{L}_{\mathrm{origin}}
    =
    -\frac{1}{N_{\mathrm{ECal}}}
    \sum_i \log p_i(y_i).
\end{equation}
The MLPF-lite objective augments this loss with soft-label
cross-entropy for the energy-fraction target,
\begin{equation}
    \mathcal{L}_{\mathrm{lite}}
    =
    \mathcal{L}_{\mathrm{origin}}
    + \lambda_f
    \left[
    -\frac{1}{N_{\mathrm{ECal}}}
    \sum_i \sum_j f_{ij}\log \widehat{f}_{ij}
    \right].
\end{equation}

For the slot model, the implemented objective is
\begin{equation}
    \mathcal{L}_{\mathrm{slot}}
    =
    \lambda_o\mathcal{L}_{\mathrm{origin}}
    +\lambda_f\mathcal{L}_{\mathrm{fraction}}
    +\lambda_s\mathcal{L}_{\mathrm{valid}}
    +\lambda_c\mathcal{L}_{\mathrm{count}},
\end{equation}
where \(\mathcal{L}_{\mathrm{valid}}\) is binary cross-entropy for the
existence of each electron slot and
\(\mathcal{L}_{\mathrm{count}}\) is categorical cross-entropy for the
event-level electron multiplicity. The slot training code supports
inverse-frequency class weights for the hit-origin and count terms, computed
from the training split.

The training entry points use the AdamW optimizer and support gradient
clipping, checkpointing and resuming. For the multi-task transformer paths,
a validation-loss-based learning-rate reduction can optionally be enabled.
Model width, number of layers, attention heads, loss coefficients, learning
rate, batch size and number of epochs are configuration parameters recorded
with each run; the values associated with reported results are therefore
specified in the Results chapter or accompanying run configuration rather
than built into the scientific definition of the task.

\subsection{Evaluation Methodology}
\label{sec:method_evaluation}

Hit-origin performance is evaluated on the held-out test split using
cross-entropy loss, hit-wise accuracy and confusion matrices. For the
multi-task fraction prediction, the pipeline additionally calculates
mean-absolute and mean-squared deviations between predicted and target
origin-energy fractions and can visualize predicted against target fractions.
These quantities distinguish incorrect dominant assignments from less severe
errors on genuinely shared hits.

For the slot model, evaluation further records per-slot validity accuracy,
exact slot-set accuracy, direct electron-count accuracy and an event-count
confusion matrix. Count accuracy can be separated by the true event
multiplicity so that two-electron and three-electron behaviour are not hidden
by a single aggregate number. The implemented visualization tools produce
training and validation loss and accuracy curves, normalized confusion
matrices, and paired three-dimensional displays of truth and predicted
\gls{ecal} hit assignments for selected test events. This chapter defines
these diagnostics; their numerical outcomes are presented in the Results
chapter.
